\begin{resumo}[Abstract]
 \begin{otherlanguage*}{english}
   Software product quality is essential to ensure maintainability and reduce costs 
   throughout the lifecycle. Quality models such as ISO/IEC 25010 highlight 
   maintainability as a fundamental characteristic, composed of attributes such as 
   analyzability, modifiability, and testability. Code smells, indicators of design and 
   implementation problems, are associated with increased maintenance effort and 
   degradation of internal quality. Traditional static analysis tools, such as SonarQube, 
   are widely used to detect these problems, but have limitations related to false 
   positives and dependence on fixed metric thresholds. Recently, Artificial Intelligence 
   (AI) techniques, especially Large Language Models (LLMs), have demonstrated potential 
   to capture more complex structural and semantic patterns. This work empirically 
   investigates how AI-based systems support software quality assurance in detecting 
   anomalies in source code. To this end, a structured literature review was conducted, 
   adapted from the systematic review protocol, which identified 42 relevant studies 
   published between 2022 and 2025. The analysis showed that, although static analysis 
   tools remain predominant in the industry, AI-based approaches have the capacity to 
   identify patterns that go beyond fixed syntactic rules. A controlled experiment was 
   then planned to compare the LLMs with SonarQube in detecting code 
   smells, using manually annotated datasets as reference. The expected results aim to 
   produce empirical evidence on the relative contribution of each approach, identifying 
   limitations and opportunities for integration between traditional methods and emerging 
   AI techniques in code quality assurance.

   \vspace{\onelineskip}
 
   \noindent 
   \textbf{Key-words}: software engineering. artificial intelligence. code smells. software quality. language models. static analysis.
 \end{otherlanguage*}
\end{resumo}
